\documentclass[11pt]{article}
\usepackage[margin=1in]{geometry}
\usepackage{amsmath,amssymb,amsthm,mathtools}
\usepackage{algorithm}
\usepackage[noend]{algpseudocode}
\usepackage{booktabs,array,enumitem}
\usepackage{xcolor}
\usepackage[colorlinks=true,linkcolor=blue!50!black,citecolor=blue!50!black,urlcolor=blue!50!black]{hyperref}
\usepackage[capitalise,nameinlink]{cleveref}

\newtheorem{theorem}{Theorem}[section]
\newtheorem{proposition}[theorem]{Proposition}
\newtheorem{lemma}[theorem]{Lemma}
\newtheorem{corollary}[theorem]{Corollary}
\theoremstyle{definition}
\newtheorem{definition}[theorem]{Definition}
\newtheorem{remark}[theorem]{Remark}
\newtheorem{example}[theorem]{Example}

\newcommand{\M}{\mathcal{M}}
\newcommand{\one}{\mathbf{1}}
\newcommand{\E}{\mathbb{E}}
\newcommand{\Zh}{\widehat{Z}}
\newcommand{\wt}{\tilde w}
\newcommand{\mask}{\texttt{m}}
\newcommand{\op}[1]{\textup{\textsc{#1}}}
\newcommand{\status}[1]{\textsf{\small[#1]}}
\newcommand{\proofref}[1]{\par\noindent\textit{Proof.} #1\par\medskip}

\title{Trajectory-exact constrained decoding for masked diffusion models:\\ routes, twists and ambiguity\\[4pt] \large Working notes}
\author{}
\date{October 2026}

\begin{document}
\maketitle

\begin{abstract}
These notes describe a decoder that samples a masked diffusion language model conditioned on a hard constraint,
in the sense of drawing from the model's own output distribution restricted to the constraint. The constraint can
be a regular expression, a counting condition, a copy relation, a context-free grammar (ambiguous grammars
included) or a tree structure. The decoder has two layers. The first is a \emph{route}: a tractable structure that
computes the probability of the constraint under a product distribution and samples from that distribution
restricted to the constraint. The second is twisted sequential Monte Carlo, which corrects the error that per-step
constrained decoding makes when the denoiser revises its predictions after seeing its own outputs. We collect the
mathematical results and algorithms behind both layers and the open directions. Results taken from the
literature are stated with a reference to their proofs. Results that are new here, or new combinations of
known pieces, are proved. Experimental results are not included.
\end{abstract}

\tableofcontents

%======================================================================
\section{Introduction}\label{sec:intro}
%======================================================================

\subsection{Masked diffusion sampling}
A masked diffusion language model (MDLM) \cite{sahoo2024} produces a sequence of length $L$ by starting from a
sequence of mask symbols and filling in positions over $T$ steps. At each step a neural \emph{denoiser} looks at
the current partial sequence and outputs, for every masked position, a categorical distribution over the
vocabulary. A \emph{scheduler} picks which positions to reveal at this step, and those positions are filled by
sampling from the denoiser's distributions, independently across positions. The revealed positions stay fixed and
the process repeats until nothing is masked. We call this procedure the \emph{native sampler}, and write $p$ for the
distribution it induces over complete sequences (and over whole trajectories, when that matters).

Two features of this sampler shape everything below. Within a step, the denoiser's predictions are a product
distribution over positions. Across steps, those predictions change, because each call of the denoiser conditions
on everything revealed so far.

\subsection{What conditional sampling asks for}
Let $C$ be a set of acceptable outputs: strings matching a pattern, balanced bracket sequences, RNA sequences that
can fold into a valid secondary structure, and so on. The goal is to sample from
\[
  p(x\mid x\in C)=\frac{p(x)\,\one_C(x)}{p(C)},
\]
the model's own distribution restricted to $C$. Rejection sampling (run the native sampler until the output lies
in $C$) achieves this exactly but is useless when $p(C)$ is small. Two cheaper methods guarantee $x\in C$.
Each changes the distribution.

\emph{Token masking} forbids, at each step, tokens that would make $C$ unreachable, and renormalises over the
rest. It ignores how much of the model's probability stays inside $C$ after each choice, so it over-represents
sequences whose early tokens leave little valid probability mass.

\emph{Step-exact decoding} \cite{de26} goes further. At each step it computes the product distribution given by the
denoiser's current predictions, restricts it to $C$, samples a complete sequence from that restricted
distribution, and keeps the positions the scheduler chose to reveal. Each step is then exact for the product
distribution the denoiser predicts at that moment. The composition over steps is still biased, as the following
example shows.

\begin{example}[Trajectory bias with two positions]\label{ex:bias}
Let $L=2$, $T=2$, reveal position 1 at the first step and position 2 at the second, and let
$C=\{(a,a):a\in V\}$, i.e.\ the two tokens must agree. Write $p^1(a)$ for the denoiser's prediction at position 1
from the all-mask state, $p^2(a\mid\varnothing)$ for its prediction at position 2 from the same state, and
$p^2(a\mid a)$ for its prediction at position 2 after position 1 has been revealed as $a$. The native conditional is
\[
  p(x=(a,a)\mid C)\ \propto\ p^1(a)\,p^2(a\mid a).
\]
Step-exact decoding draws position 1 from the restricted product at the first step, which gives $a$ with
probability $\propto p^1(a)\,p^2(a\mid\varnothing)$; the second step is then forced. The two laws agree only if
$p^2(a\mid a)/p^2(a\mid\varnothing)$ does not depend on $a$, i.e.\ only if revealing position 1 does not change
the relative predictions at position 2. For a real language model it almost always does.
\end{example}

The ratio $p^2(a\mid a)/p^2(a\mid\varnothing)$ in \cref{ex:bias} is exactly the importance weight that corrects
the error. It compares the probability of the constraint under the \emph{updated} predictions with its
probability under the \emph{old} predictions, for the same partial sequence. TWISTER \cite{twister} applies
this correction at every step and resamples a population of particles, using twisted sequential Monte Carlo.
It proves that the result targets the native conditional exactly, for constraints given by deterministic finite
automata (DFAs).

\subsection{The approach in these notes}
We separate the two concerns.
\begin{enumerate}[nosep]
  \item A \emph{route} handles everything that happens within one step. Given per-position weights, it computes
    the probability of $C$ under the product distribution, samples from the product distribution restricted to
    $C$, and checks membership in $C$. Automata are one kind of route. Grammar charts, factor graphs and
    determinants are others, and each covers constraint families the others cannot handle efficiently.
  \item A \emph{trajectory layer} handles what happens across steps. It is TWISTER's twisted SMC with the
    automaton replaced by any route.
\end{enumerate}
The main theoretical observation (\cref{thm:traj}) is that the trajectory layer needs very little from the
route. The route's sampler must draw from the same measure whose total its mass computation returns, and the last
step must check membership in $C$ exactly. The route does not have to be unambiguous. An ambiguous grammar, which
counts a string once per parse tree, gives the right answer as the number of particles grows. Ambiguity only
increases the variance.

The other main result concerns resampling (\cref{sec:resampling}). The SMC estimate of $p(C)$ is unbiased for
every number of particles. When the denoiser's predictions at some intermediate step badly understate how likely
the constraint is to be met in the end, adaptive resampling discards particles that would later have recovered.
A non-trivial fraction of the estimate's expectation can then sit in events too rare to observe. We characterise
this hidden mass exactly and give two remedies that keep the estimate unbiased and have guarantees that do not
depend on the quality of the predictions.

\subsection{What is known and what is new}
Much of the step-level theory is known in the probabilistic-circuits literature, often in more general algebraic
form. \Cref{sec:step} states it in our notation with references. The route-generic trajectory correction, the
analysis of resampling, the remedies, the ambiguity bound and the superset construction are new here, as far as
we know. Each result is tagged \status{background} (known, cited, proof referenced), \status{ours} (proved here)
or \status{proposed} (a direction, not a result). \Cref{sec:status} lists all of them.

\subsection{Outline}
\Cref{sec:smc} reviews the parts of sequential Monte Carlo we use. \Cref{sec:setup} fixes notation.
\Cref{sec:step} gives the step-level theory and \cref{sec:routes} the routes and their samplers.
\Cref{sec:traj} contains the trajectory-exact decoder and its correctness theorem. \Cref{sec:resampling}
analyses resampling, \cref{sec:ambiguity} ambiguity, and \cref{sec:beyond} constraints that need more than an exact
route. \Cref{sec:approach} lists open directions.

%======================================================================
\section{Background: twisted sequential Monte Carlo}\label{sec:smc}
%======================================================================

This section collects standard facts; \cite{delmoral2004} is the reference for all of them.

\paragraph{Importance sampling on paths.}
Let $\gamma$ be an unnormalised target on trajectories $x_{0:T}$ with normaliser $\mathcal Z=\sum\gamma$, and let
$q$ be a proposal with $q>0$ wherever $\gamma>0$. The weight $W=\gamma/q$ satisfies $\E_q[W]=\mathcal Z$, and the
weighted samples $\{(x^k,W^k)\}$ estimate expectations under $\gamma/\mathcal Z$. Here
$\gamma(x_{0:T})=p(x_{0:T})\one_C(x_T)$, so $\mathcal Z=p(C)$, which we write $h_0$.

\paragraph{Sequential form and twists.}
If $p$ and $q$ factorise over steps, $p=\prod_t\kappa_{t+1}$ and $q=\prod_tM_{t+1}$, the weight can be accumulated
step by step. A \emph{twist} is a sequence of non-negative functions $\psi_0,\dots,\psi_T$ with $\psi_T=\one_C$.
It defines intermediate targets $\gamma_t(x_{0:t})=p(x_{0:t})\psi_t(x_t)$ and \emph{potentials}
\begin{equation}\label{eq:G}
  G_{t+1}(x_t,x_{t+1})=\frac{\kappa_{t+1}(x_{t+1}\mid x_t)\,\psi_{t+1}(x_{t+1})}{M_{t+1}(x_{t+1}\mid x_t)\,\psi_t(x_t)}.
\end{equation}
The product $\psi_0(x_0)\prod_tG_{t+1}$ telescopes to $\gamma/q$, so the final weight is unchanged by the choice of
twist. The twist only changes the intermediate weights, and therefore which particles survive resampling. The ideal
twist is the \emph{Doob function} $h_t(x)=\Pr_p(X_T\in C\mid X_t=x)$. With $\psi_t=h_t$, the intermediate weights
reflect exactly how likely each particle is to end in $C$.

\paragraph{Resampling and the normalising-constant estimate.}
With $K$ particles, the weights tend to concentrate on a few of them. A common diagnostic is the effective sample
size $\mathrm{ESS}=(\sum_kW^k)^2/\sum_k(W^k)^2\in[1,K]$. Adaptive resampling draws $K$ new particles from the
current ones in proportion to their weights whenever $\mathrm{ESS}$ falls below a threshold $\tau K$, and then
resets the weights. The estimate
\[
  \Zh=\prod_{\text{resampling times}}\Big(\frac1K\sum_kW^k\Big)\times\Big(\frac1K\sum_kW^k\Big)_{\text{final}}
\]
is unbiased for $\mathcal Z$ at every $K$ \cite{delmoral2004}. The weighted particles are consistent for the
normalised target as $K\to\infty$.

%======================================================================
\section{Setup and notation}\label{sec:setup}
%======================================================================

\paragraph{Tokens and classes.}
Let $V$ be the vocabulary and $\mask\notin V$ the mask symbol. Constraints are usually stated over a small alphabet
$\Sigma$ of token \emph{classes}, through a map $\phi:V\to\Sigma$ (digit, opening bracket, nucleotide). We write
$C\subseteq\Sigma^L$, and a token sequence $x$ satisfies $C$ if $\phi(x)\in C$. For a partial sequence
$x\in(V\cup\{\mask\})^L$, $\M(x)$ is the set of its masked positions.

\paragraph{Native kernel.}
Following \cite{twister}, one native step from $x$ to $x'$ has probability
\begin{equation}\label{eq:native}
  \kappa_{t+1}(x'\mid x)=s_{t+1}(r\mid x)\prod_{i\in r}p^i_\theta(x'_i\mid x),
\end{equation}
where $r\subseteq\M(x)$ is the set of positions revealed at the step, $s_{t+1}$ is the scheduler (for example,
$L/T$ positions chosen uniformly among the masked ones) and $p^i_\theta(\cdot\mid x)$ is the denoiser's prediction at
position $i$. The native path law is $p(x_{0:T})=\prod_t\kappa_{t+1}(x_{t+1}\mid x_t)$, with $x_0=\mask^L$ and
$x_T$ fully revealed.

\paragraph{Target.}
The \emph{Doob path law} $p^\star(x_{0:T})\propto p(x_{0:T})\one_C(x_T)$ has terminal marginal $p(\cdot\mid C)$.
It is a Markov chain with kernel $\kappa_{t+1}h_{t+1}/h_t$. The Doob function $h_t$ depends on the denoiser's
outputs at every future state, so it cannot be computed directly.

\paragraph{Frozen weights and the clamped sum.}
For a partial sequence $x$, the class weights are
\[
  w_i(\sigma\mid x)=\sum_{v:\,\phi(v)=\sigma}p^i_\theta(v\mid x)\qquad(i\in\M(x)),
\]
and revealed positions are clamped to their class, $w_i(\sigma\mid x)=[\sigma=\phi(x_i)]$. The \emph{clamped sum} of
a partial sequence $x'$ under the predictions made at $x$ is
\begin{equation}\label{eq:clamped}
  Z(x';x)=\sum_{y\in C}\ \prod_{i\in\M(x')}w_i(y_i\mid x)\prod_{i\notin\M(x')}[y_i=\phi(x'_i)].
\end{equation}
$\eta_t(x):=Z(x;x)$ is the probability that $C$ holds if every masked position is filled independently from the
current predictions. TWISTER uses $\eta_t$ in place of $h_t$. We call it the \emph{frozen twist}, because it
treats the current predictions as if they would not change.

%======================================================================
\section{Step-level theory}\label{sec:step}
%======================================================================

\subsection{The constraint polynomial}
Every per-step quantity in constrained decoding is a query on a single polynomial.

\begin{definition}[Constraint polynomial]
For $C\subseteq\Sigma^L$ and $w\in\mathbb{R}_{\ge0}^{L\times\Sigma}$,
$P_C(w)=\sum_{y\in C}\prod_{i=1}^Lw_{i,y_i}$.
\end{definition}

The frozen twist is $\eta_t(x)=P_C(w(x))$. The constrained marginals are
$\mu_{i,\sigma}=\Pr_q(y_i=\sigma\mid C)=w_{i,\sigma}\,\partial_{i,\sigma}P_C(w)/P_C(w)$, and step-exact decoding needs
samples from $q(y\mid C)\propto\prod_iw_{i,y_i}\one_C(y)$. $P_C$ is set-multilinear (each monomial contains one
variable per position), so the marginal formula holds for any circuit that computes $P_C$.

\begin{definition}[Structure circuit and derivations]
A \emph{structure circuit} is a smooth, decomposable, monotone arithmetic circuit over the leaves $w_{i,\sigma}$
with unit parameters. Automaton trellises, inside charts of grammars restricted to length $L$, and variable
elimination with $0/1$ factors are all of this kind. A structure circuit $S$ computes
\[
  P_S(w)=\sum_{y\in\Sigma^L}d_S(y)\prod_iw_{i,y_i},
\]
where $d_S(y)\in\mathbb N$ is the number of \emph{derivations} of $y$ in $S$ (accepting paths of an automaton,
parse trees of a grammar). $S$ represents $C$ if $\operatorname{supp}d_S=C$, and is \emph{unambiguous} if
$d_S\equiv1$ on $C$.
\end{definition}

\subsection{Semirings and ambiguity}
\begin{proposition}[\status{background}]\label{prop:semiring}
Let $S$ represent $C$. Evaluated on leaves $w\ge0$:
\begin{center}
\begin{tabular}{lll}
\toprule
semiring & value of $S$ & equals the constrained query\\
\midrule
Boolean $(\lor,\land)$ & $[\exists y\in C$ consistent with $w]$ & always (completability)\\
max-times & $\max_{y\in C}\prod_iw_{i,y_i}$ & always (MAP, recovered by backtracking)\\
sum-product & $\sum_yd_S(y)\prod_iw_{i,y_i}$ & if and only if $d_S\equiv1$ on $C$\\
\bottomrule
\end{tabular}
\end{center}
\end{proposition}
\proofref{A special case of Wang et al.\ \cite{wang2024atlas}, Thm.~4 and Ex.~2: mapping a Boolean circuit into a
semiring is tractable when the semiring is idempotent or the circuit is deterministic.}

Every derivation of $y$ has weight $\prod_iw_{i,y_i}$. Idempotent addition counts the $d_S(y)$ copies once, and
ordinary addition counts all of them. Top-down sampling from a smooth, decomposable, monotone circuit is correct
for the polynomial the circuit computes, deterministic or not. Ambiguity therefore makes a sampler draw from the
derivation-weighted distribution $\propto d_S(y)\prod_iw_{i,y_i}$. Dang and Ermon call their NFA variant a ``soft
proxy'' for this reason \cite{de26}.

\begin{corollary}[Exact MAP for any CFG \status{background}]
For every context-free grammar $G$, including inherently ambiguous ones, Viterbi on the length-$L$ chart computes
$\arg\max_{y\in L(G)\cap\Sigma^L}\prod_iw_{i,y_i}$ in $O(L^3|G|)$.
\end{corollary}
\proofref{The max-times row of \cref{prop:semiring} applied to the chart of \cref{sec:chart}.}

\subsection{When is the step-level problem tractable?}
\begin{theorem}[\status{background}]\label{thm:char}
Fix $\Sigma$ and a family $\{C_L\}$. With all bounds polynomial in $L$ and $|\Sigma|$, the following are equivalent.
\begin{enumerate}[label=(\alph*),nosep]
  \item $P_{C_L}$ has a polynomial-size arithmetic circuit over $\mathbb R$. Subtraction is allowed; the circuit
    need not be decomposable, deterministic or monotone.
  \item A real-RAM algorithm computes $P_{C_L}(w)$ for every $w\ge0$.
  \item A real-RAM algorithm computes all constrained marginals $\mu_{i,\sigma}(w)$ whenever $P_{C_L}(w)>0$.
\end{enumerate}
Given a circuit of size $s$, $P_C(w)$ costs $O(s)$, all $L|\Sigma|$ marginals together cost $O(s)$, and an exact
sample from $q(\cdot\mid C)$ costs $O(Ls)$. If the circuit is smooth, decomposable and monotone, one top-down pass
gives a sample.
\end{theorem}
\proofref{(b)$\Rightarrow$(a) is the completeness theorem of Broadrick et al.\ \cite{broadrick2025}, Prop.~5.4
(Sec.~5.1, proof in App.~C). Their virtual-evidence marginalisation (Def.~5.1) is query (b) for $\one_C$ after
one-hot encoding each position. Their proof uses Strassen's division elimination and a polynomial-identity
argument on an open box; we checked that both steps go through for $P_C$ with correctness only on the positive
orthant. (a)$\Rightarrow$(b) is evaluation, and all marginals come from one reverse-mode differentiation pass.
(c)$\Rightarrow$(b) and the $O(Ls)$ sampler are self-reducibility under clamping \cite{jvv86}.}

Broadrick et al.\ also show (their Prop.~4.9 and Thm.~4.10) that tractable marginals under hard evidence alone do
not imply a small circuit. Condition (b) needs arbitrary weights, which decoding supplies, since the weights are
the denoiser's predictions.

\begin{remark}[Approximate sampling \status{background}]
A polynomial-time almost-uniform sampler for $q(\cdot\mid C)$ that works for every clamped $w$ yields an FPRAS
for $P_C$ \cite{jvv86}. An efficient sampler therefore implies
efficient \emph{approximate} counting, which is weaker than (a) to (c). The route compiler uses the converse to
decide when to refuse a constraint (\cref{sec:beyond}).
\end{remark}

\begin{remark}[Decomposable circuits do not cover everything \status{background}]\label{rem:arb}
Encode a tree on nodes $\{0,\dots,n\}$ by head pointers $y_m\in\{0,\dots,n\}\setminus\{m\}$, and let $C$ be the
set of arborescences rooted at $0$. Then $P_C$ is the directed spanning-tree polynomial. It has a polynomial-size
circuit via the directed matrix-tree theorem (\cref{sec:det}). Every monotone circuit for it has size $2^{\Omega(n)}$
\cite{jerrumsnir1982}. Smooth decomposable circuits with non-negative parameters are monotone, so no automaton,
grammar or bounded-treewidth route is polynomial for this constraint, while exact decoding through a determinant
is.
\end{remark}


%======================================================================
\section{Routes and their exact samplers}\label{sec:routes}
%======================================================================

\begin{definition}[Route]
A \emph{route} for $C$ provides three operations on weights $w\ge0$:
\begin{itemize}[nosep]
  \item $\op{Mass}(w)=\sum_yd(y)\prod_iw_{i,y_i}$, for a fixed $d\ge0$ with $\operatorname{supp}d\supseteq C$;
  \item $\op{Sample}(w)$, an exact draw from $y\propto d(y)\prod_iw_{i,y_i}$;
  \item $\op{Accepts}(y)=\one_C(y)$, an exact membership test.
\end{itemize}
The route is \emph{exact} if $d=\one_C$, \emph{derivation-weighted} if $d\in\mathbb N$ and $\operatorname{supp}d=C$,
and a \emph{superset route} if $\operatorname{supp}d\supsetneq C$.
\end{definition}

The condition the trajectory layer relies on is that $\op{Sample}$ draws from exactly the measure whose total
$\op{Mass}$ returns. The routes below all satisfy it. In each case the correctness of the mass computation and
the sampler is standard; we give the recursions and references.

\subsection{Automata}
For a finite automaton $(Q,q_0,\delta,F)$ over $\Sigma$, the forward recursion
\[
  \alpha_0(q)=[q=q_0],\qquad\alpha_i(q)=\sum_{q',\sigma:\,q\in\delta(q',\sigma)}\alpha_{i-1}(q')\,w_{i,\sigma},
  \qquad\op{Mass}=\sum_{q\in F}\alpha_L(q),
\]
costs $O(L|\delta|)$. $\op{Sample}$ is forward filtering, backward sampling. Draw
$q_L\propto\alpha_L(q)[q\in F]$; then, for $i=L,\dots,1$, draw $(q_{i-1},\sigma_i)$ with probability
$\propto\alpha_{i-1}(q_{i-1})w_{i,\sigma_i}[q_i\in\delta(q_{i-1},\sigma_i)]$. For a DFA, $d=\one_C$. For an NFA,
$d(y)$ is the number of accepting paths of $y$.

Constraints that compile to automata include forbidden and required motifs (via an Aho--Corasick
automaton), counts such as ``exactly $k$ tokens of class $c$'' or ``a fraction within a band'' (a
counter automaton with $O(Lk)$ states in any variable order), regular expressions over classes, bounded-depth
bracketing, sliding-window all-different constraints, and products of these.

\subsection{Grammar charts}\label{sec:chart}
For a grammar in Chomsky normal form with nonterminals $N$, binary rules $A\to BC$ and lexical rules $A\to\sigma$,
the inside algorithm on the chart of length $L$,
\begin{gather*}
  \beta(A,i,i{+}1)=\sum_\sigma[A\to\sigma]\,w_{i,\sigma},\qquad
  \beta(A,i,j)=\sum_{A\to BC}\ \sum_{i<k<j}\beta(B,i,k)\,\beta(C,k,j),\\
  \op{Mass}=\beta(S,0,L),
\end{gather*}
costs $O(L^3|G|)$. The leaves for all positions come from one matrix product of the $L\times|V|$ denoiser output
with the $|V|\times|N_{\mathrm{pre}}|$ lexical map.

\begin{algorithm}[h]
\caption{Top-down chart sampler}\label{alg:chart}
\begin{algorithmic}[1]
\State push $(S,0,L)$
\While{the stack is non-empty}
  \State pop $(A,i,j)$
  \If{$j=i+1$} draw $\sigma\propto[A\to\sigma]\,w_{i,\sigma}$ and set $y_i\gets\sigma$
  \Else\ draw $(A\to BC,\,k)\propto\beta(B,i,k)\,\beta(C,k,j)$; push $(B,i,k)$ and $(C,k,j)$
  \EndIf
\EndWhile
\end{algorithmic}
\end{algorithm}

\Cref{alg:chart} draws each parse tree with probability proportional to its weight. The
sampled string therefore has probability $\propto d(y)\prod_iw_{i,y_i}$, where $d(y)$ is its number of parse
trees. On a GPU, all spans of one
length are processed together and the draws are batched over particles.

\begin{example}[Grammars used in the experiments]
\leavevmode
\begin{itemize}[nosep]
  \item Unambiguous Dyck-1, $S\to(S)S\mid\varepsilon$, converted to CNF.
  \item Ambiguous Dyck-1, $S\to SS\mid(S)\mid()$. It generates the same language, but a string with several
    adjacent blocks has several parse trees.
  \item An RNA pairing grammar in the style of Nussinov and Jacobson \cite{nussinov1980}. Every base is paired,
    pairs do not cross, and each pair is one of AU, UA, GC, CG, GU, UG. Parse trees correspond one to one with
    secondary structures, so $d(y)$ is the number of structures $y$ can fold into. This number can grow
    exponentially with $L$; $(GC)^n$ has Catalan$(n)$ structures.
\end{itemize}
\end{example}

\subsection{Factor graphs}
If $C$ is given by $0/1$ factors and an elimination order has induced width $\omega$, variable elimination
computes $\op{Mass}$ in $O(L|\Sigma|^{\omega+1})$. $\op{Sample}$ eliminates in that order and
samples in reverse. Derivations here are assignments, so $d=\one_C$. The case that motivates this route is the
same-order copy relation $y_i=y_{i+k}$ for $i\le k$ \cite{sz26}, which has width 1:
\[
  \op{Mass}=\prod_{i\le k}\ \sum_\sigma w_{i,\sigma}w_{i+k,\sigma}\times\prod_{i>2k}\ \sum_\sigma w_{i,\sigma}.
\]
An automaton that reads positions in order needs $|\Sigma|^k$ states for the same relation.

\subsection{Determinants}\label{sec:det}
Let $w_m(h)$ be the weight that node $m$ has head $h$. Let $\hat L$ be the $n\times n$ matrix on the non-root
nodes with $\hat L_{mm}=\sum_{h\ne m}w_m(h)$ (the root included in the sum) and $\hat L_{hm}=-w_m(h)$ for $h\ne m$.
The directed matrix-tree theorem gives $\op{Mass}=\det\hat L$. Differentiating $\log\det\hat L$
gives the marginals
\[
  \mu_m(h)=w_m(h)\big([\hat L^{-1}]_{mm}-[h\ne0]\,[\hat L^{-1}]_{mh}\big)
\]
(see \cite{koo2007}). Both cost $O(n^3)$. $\op{Sample}$ is Wilson's algorithm \cite{wilson1996}: for each node not
yet in the tree, run the random walk that moves from $m$ to $h$ with probability $w_m(h)/\sum_{h'}w_m(h')$ until
it reaches the tree, erase the loops of the path and add the path to the tree. The resulting arborescence has
probability $\propto\prod_mw_m(y_m)$. By \cref{rem:arb} no polynomial monotone route exists for this constraint.

\subsection{From classes to tokens}
A route draws a class sequence $y$. Each revealed position then draws a token within its class,
$v\propto p^i_\theta(v\mid x)[\phi(v)=y_i]$. The combined token-level law is
$\prod_ip^i_\theta(v_i\mid x)\one_C(\phi(v))/Z$, which is what the trajectory layer requires.

\subsection{Route compiler}
The compiler turns a constraint specification into a route. It estimates the cost of one pass for each candidate
from compile-time statistics ($O(L|\delta|)$ for a chain, $O(L^3|G|)$ for a chart, $O(L|\Sigma|^{\omega+1})$ for
variable elimination, $O(n^3)$ for a determinant) and picks the cheapest. Which route is cheapest depends on the
representation and the variable order as well as on the constraint: copy is exponential as an in-order automaton
and linear as a factor graph. When no candidate fits the budget, or the family has \#P-hard counting (for
example, CNF clauses over positions), the compiler refuses instead of approximating.

%======================================================================
\section{Trajectory-exact decoding with any route}\label{sec:traj}
%======================================================================

\subsection{The step-exact proposal}
Write $Z_d(x';x)$ for the clamped sum \eqref{eq:clamped} with $\one_C$ replaced by the route's $d$, and
$\eta_t(x)=Z_d(x;x)$.

\begin{lemma}[Step-exact proposal \status{background}]\label{lem:prop}
The kernel
\[
  M_{t+1}(x'\mid x)=s_{t+1}(r\mid x)\prod_{i\in r}p^i_\theta(x'_i\mid x)\,\frac{Z_d(x';x)}{Z_d(x;x)}
\]
is a probability kernel. It is sampled by drawing $r$ from the native scheduler, drawing a complete class sequence
$y\sim\op{Sample}(w(x))$ with revealed positions clamped, drawing tokens within classes at the positions in $r$, and
revealing only those positions.
\end{lemma}
\proofref{TWISTER \cite{twister}, Prop.~3.6, proof in App.~A, stated for DFAs. The proof uses only the fact that
the sampler draws from the measure whose total is $Z_d(x;x)$, so it applies to any route.}

Combining \cref{lem:prop} with \eqref{eq:G}, the scheduler and the token probabilities cancel, and the potential
becomes
\begin{equation}\label{eq:potential}
  G_{t+1}(x,x')=\frac{Z_d(x;x)}{\psi_t(x)}\cdot\frac{\psi_{t+1}(x')}{Z_d(x';x)}.
\end{equation}
TWISTER takes $\psi_t=\eta_t$ for $t<T$, which gives $G_{t+1}=Z_d(x';x')/Z_d(x';x)$: the mass of the new partial
sequence under the new predictions, divided by its mass under the old predictions. In \cref{ex:bias} this ratio
is $p^2(a\mid a)/p^2(a\mid\varnothing)$.

\subsection{The algorithm}
\Cref{alg:smc} is TWISTER with the automaton replaced by a route. It has two optional parameters, introduced in
\cref{sec:resampling}: defensive mixing $\epsilon$ and twist power $\beta$. With $\epsilon=0$ and $\beta=1$ it is
TWISTER's algorithm. With $K=1$ and the weights ignored, it is step-exact decoding \cite{de26}.

\begin{algorithm}[h]
\caption{Route-generic twisted SMC (one run with $K$ particles)}\label{alg:smc}
\begin{algorithmic}[1]
\Require route $(\op{Mass},\op{Sample},\op{Accepts})$; denoiser; $T$; $K$; ESS threshold $\tau$;
  $\epsilon\in[0,1)$; $\beta\in(0,1]$
\State for each $k$: $x^k\gets\mask^L$, $\eta^k\gets\op{Mass}(w(x^k))$, $\log W^k\gets\beta\log\eta^k$; set $\log\Zh\gets0$
\For{$t=0,\dots,T-1$}
  \For{each particle $k$}
    \State draw $r\sim s_{t+1}(\cdot\mid x^k)$ and $y\sim\op{Sample}(w(x^k))$; draw tokens within classes on $r$;
      $x'\gets x^k$ revealed on $r$
    \State $z\gets Z_d(x';x^k)$ \Comment{mass of $x'$ under the old predictions}
    \If{$t<T-1$} call the denoiser at $x'$; $\eta'\gets Z_d(x';x')$; $\psi'\gets(\eta')^\beta$
    \Else\ $\psi'\gets\op{Accepts}(x')$ \Comment{exact membership at the last step}
    \EndIf
    \State $\log W^k\gets\log W^k+(1-\beta)\log\eta^k+\log\psi'-\log z$; \ $x^k\gets x'$; \ $\eta^k\gets\eta'$
  \EndFor
  \If{$t<T-1$ and $\mathrm{ESS}(W)<\tau K$}
    \State $\log\Zh\gets\log\Zh+\log\frac1K\sum_kW^k$; \ $\wt\gets W/\sum_kW^k$; \ $\pi\gets(1-\epsilon)\wt+\epsilon/K$
    \State draw ancestors $a_1,\dots,a_K$ from $\pi$ (independently, or systematically); for each $k$:
      $x^k\gets x^{a_k}$, $\eta^k\gets\eta^{a_k}$, $W^k\gets\wt_{a_k}/\pi_{a_k}$
  \EndIf
\EndFor
\State $\log\Zh\gets\log\Zh+\log\frac1K\sum_kW^k$
\State \Return weighted particles $\{(x^k,W^k)\}$ and $\Zh$
\end{algorithmic}
\end{algorithm}

Per particle and step, the algorithm makes one denoiser call, one route sample and two mass evaluations: the
clamped pass under the old predictions and the pass under the new ones. The latter is reused at the next step.

\subsection{Correctness}
\begin{theorem}[Route-generic trajectory exactness \status{ours}]\label{thm:traj}
Assume the denoiser gives every token positive probability. Let the route satisfy $\operatorname{supp}d\supseteq
C$, with ambiguity allowed. Let the twists satisfy $\psi_t(x)>0$ whenever $h_t(x)>0$, and $\psi_T=\one_C$; for
instance $\psi_t=\eta_t^\beta$ with any $\beta>0$. Let $W(x_{0:T})=\psi_0(x_0)\prod_{t<T}G_{t+1}$ and
$q=\prod_tM_{t+1}$. Then:
\begin{enumerate}[label=(\roman*),nosep]
  \item $q(x_{0:T})\,W(x_{0:T})=p(x_{0:T})\,\one_C(x_T)$ for every path;
  \item the weighted particles of \cref{alg:smc} target the Doob path law $p^\star$ and are consistent for
    $p(\cdot\mid C)$ as $K\to\infty$;
  \item $\E[\Zh]=h_0$ for every $K\ge1$, every $\epsilon\in[0,1)$, and every resampling rule whose decision at
    each step depends only on the current particles and weights;
  \item the derivation count affects only the variance: the factor $d(x_T)$ introduced by the last proposal step
    cancels in the last potential.
\end{enumerate}
\end{theorem}
\begin{proof}
(i) By \cref{lem:prop}, $q=\prod_t\kappa_{t+1}Z_d(x_{t+1};x_t)/Z_d(x_t;x_t)$. Multiplying by
\eqref{eq:potential}, the factors $Z_d$ cancel, and the twists telescope to
$qW=\psi_0\prod_t\kappa_{t+1}\psi_{t+1}/\psi_t=p\,\psi_T=p\,\one_C$. The divisions are legitimate on the relevant
paths. If $x'$ can still be completed to an element of $C$, some $y\in C$ agrees with $x'$; then $d(y)>0$, so
$Z_d(x';x)>0$ under the positivity assumption, and $M_{t+1}(x'\mid x)>0$. If $\psi_t(x)=0$, the particle already
has weight zero, and we set $0/0=0$.

(ii) follows from (i) by the standard Feynman--Kac argument \cite{delmoral2004}.

(iii) Unbiasedness of $\Zh$ for multinomial resampling with adaptive decisions is standard \cite{delmoral2004}.
For $\epsilon>0$ we check that the reweighted resampling step preserves conditional expectations. Let
$\mathcal F$ be the particle system just before a resampling step, with weights $W^k$ and mean $\bar W$. For any
$f$,
\[
  \E\Big[\bar W\,\frac1K\sum_j\frac{\wt_{a_j}}{\pi_{a_j}}f(x^{a_j})\Bigm|\mathcal F\Big]=\bar W\sum_k\wt_kf(x^k)
  =\frac1K\sum_kW^kf(x^k),
\]
because $\pi_k>0$ wherever $\wt_k>0$. A resampling step therefore leaves the conditional expectation of the
weighted particle measure unchanged, whether or not it fires. Iterating over steps gives $\E[\Zh]=\E_q[W]$, which
equals $h_0$ by (i).

(iv) At the last step $x_T$ is fully revealed, so $Z_d(x_T;x_{T-1})=d(x_T)$, and \eqref{eq:potential} gives
$G_T=Z_d(x_{T-1};x_{T-1})\one_C(x_T)/\big(\psi_{T-1}(x_{T-1})\,d(x_T)\big)$.
\end{proof}

The theorem separates what each kind of decoding needs. Step-exact decoding needs an exact route
(\cref{thm:char}). Trajectory-exact decoding needs a route whose sampler matches its mass, a twist that is
positive on completable states, and an exact membership test at the end. Any NFA or CFG chart, ambiguous or not,
meets these conditions in polynomial time per particle. This does not contradict the \#P-hardness of computing the
exact future-validity function for general CFGs, because SMC never computes $h_t$; it is exact only in the limit
$K\to\infty$.

\begin{proposition}[Terminal twist \status{ours}]
If an ambiguous route's own value $Z_d(x_T;x_T)=d(x_T)$ is used as the terminal twist in place of $\one_C$, the
particles target $\propto p(x_{0:T})\,d(x_T)$, the derivation-weighted law.
\end{proposition}
\begin{proof}
Replace $\psi_T=\one_C$ by $\psi_T=d$ in the telescoping product of \cref{thm:traj}(i).
\end{proof}

\begin{remark}[Relation to prior work]
For unambiguous routes, \cref{thm:traj} is TWISTER's argument \cite{twister} (Thm.~4.2, proof in App.~D)
applied to a different route. Their proof uses the DFA only to compute $Z$. Dang et al.\ \cite{dang2026icml}
(Thm.~4.3 and the remark that follows it) observe, for NFAs in autoregressive SMC, that the derivation-weighted
distribution is still a valid proposal. \Cref{thm:traj} combines these observations for masked diffusion and for
every route of \cref{sec:routes}. We have not found the terminal-twist pitfall or the analysis of
\cref{sec:resampling} stated elsewhere.
\end{remark}

\subsection{Choosing the twist separately from the proposal}
The proof of \cref{thm:traj} uses two facts: the proposal density can be computed, and the twists telescope. The
twist can therefore be chosen freely.

\begin{proposition}[Free intermediate twists \status{background}]\label{prop:free}
Let $M_{t+1}$ be any kernel with computable density and $M_{t+1}(x'\mid x)>0$ whenever
$\kappa_{t+1}(x'\mid x)h_{t+1}(x')>0$. Let $\psi_t\ge0$ with $\psi_t>0$ where $h_t>0$, and $\psi_T=\one_C$. With
the potentials \eqref{eq:G}, $qW=p\,\one_C$, and (ii) and (iii) of \cref{thm:traj} hold. The same is true for
random twists $\psi_t(x,\xi_t)$, where the auxiliary randomness $\xi_t$ is drawn afresh at each step and the same
realisation is used in the numerator at step $t$ and the denominator at step $t+1$.
\end{proposition}
\proofref{The deterministic case is the standard twisted-SMC construction \cite{delmoral2004}; \cite{zhao2024} uses it for
autoregressive language models. For random twists, include $\xi_t$ in the state; each realisation appears once in a
numerator and once in a denominator, so the product still telescopes path by path.}

With the step-exact proposal, the twist used to decide which particles survive can thus be chosen independently of
the route used to propose. For example, $\psi_t=\eta_t^\beta$ gives the potential
$G_{t+1}=\eta_t(x)^{1-\beta}\eta_{t+1}(x')^\beta/Z_d(x';x)$, which is the update in \cref{alg:smc}.

\subsection{A two-sided identity}
\begin{lemma}[\status{ours}]\label{lem:twosided}
Under the conditions of \cref{thm:traj},
\[
  \E_q[W]=h_0,\qquad\E_p\big[\one_C(x_T)/W(x_{0:T})\big]=\Pr\nolimits_q(x_T\in C).
\]
The second expectation is over native paths. For an exact or derivation-weighted route, $\Pr_q(x_T\in C)=1$.
\end{lemma}
\begin{proof}
By \cref{thm:traj}(i), $\E_p[\one_C/W]=\sum_{\text{paths}}p\,\one_C\,q/(p\,\one_C)$, which is the $q$-probability of
ending in $C$.
\end{proof}
The first identity is checked with proposal runs, the second with native samples. Native paths whose weight exceeds
every observed proposal weight point to regions the proposal rarely visits. Running \cref{alg:smc} with resampling switched off
alongside the main run is also useful, because it separates the effect of the proposal and twist from the effect
of resampling.

%======================================================================
\section{Resampling with an over-confident twist}\label{sec:resampling}
%======================================================================

\Cref{thm:traj}(iii) says $\E[\Zh]=h_0$. If the distribution of $\Zh$ has most of its expectation in rare
events, the average over any feasible number of runs can still fall well short of $h_0$. This section shows how
adaptive resampling with the frozen twist can create such rare events, and gives two remedies. The statements
are exact and make no assumption about the denoiser.

\subsection{Continuation values and hidden mass}
\begin{lemma}[Continuation value \status{ours}]\label{lem:cont}
Without further resampling, the expected product of the remaining potentials from state $x$ at step $s$ is
\[
  V_s(x):=\E_q\Big[\prod_{t\ge s}G_{t+1}\Bigm|x_s=x\Big]=\frac{h_s(x)}{\psi_s(x)}.
\]
\end{lemma}
\begin{proof}
From step $s$ on, the product of the proposal kernels and the potentials telescopes to $p\,\psi_T/\psi_s$ along
every path. Summing over continuations gives $\Pr_p(X_T\in C\mid X_s=x)/\psi_s(x)$.
\end{proof}

Suppose resampling fires at step $s$, with cumulative weights $W^k_s$ and normalised weights $\wt_k$. By
\cref{thm:traj}(iii) and \cref{lem:cont}, any unbiased continuation of the run has
\[
  \E[\Zh\mid\mathcal F_s]=\sum_kc_k,\qquad c_k:=\frac1KW^k_s\,\frac{h_s(x^k_s)}{\psi_s(x^k_s)}.
\]
Write $W^k_s=\rho^k_s\,\psi_s(x^k_s)$, where $\rho^k_s$ is the ratio $p/q$ along particle $k$'s path up to $s$.
Then
\begin{equation}\label{eq:mismatch}
  \wt_k\propto\rho^k_s\,\psi_s(x^k_s),\qquad c_k\propto\rho^k_s\,h_s(x^k_s).
\end{equation}
Resampling allocates offspring in proportion to $\psi_s$, while each particle's share of the expectation is
proportional to $h_s$.

\begin{proposition}[Hidden mass under multinomial resampling \status{ours}]\label{prop:hidden}
Under multinomial resampling ($\pi=\wt$), particle $k$ has at least one offspring with probability
$1-(1-\wt_k)^K\le K\wt_k$; when it does, its contribution to the run's estimate is multiplied by about
$1/(K\wt_k)$. For $\varepsilon>0$ define the hidden mass
\[
  H_\varepsilon=\frac1{h_0}\,\E\Big[\sum_kc_k\,\one\{\wt_k<\varepsilon\}\Big].
\]
Over $R$ independent runs, the expected number of runs in which some particle with $\wt_k<\varepsilon$ has an
offspring is at most $RK\varepsilon\,\E[N_\varepsilon]$, where $N_\varepsilon$ is the number of such particles in
a run. When this number is much smaller than 1, the average of $\Zh$ over the $R$ runs misses a fraction of about
$H_\varepsilon$ of $h_0$ with high probability, although $\E[\Zh]=h_0$.
\end{proposition}
\begin{proof}
The offspring count of particle $k$ is binomial with parameters $K$ and $\wt_k$, which gives the first bound; a
union bound over particles and runs gives the second. On the event that particle $k$ has $n_k$ offspring, they
carry weight $\bar W$ each, against an expected share $c_k$ that \cref{lem:cont} assigns to a weight of
$W^k_s=K\bar W\wt_k$; hence the factor $n_k/(K\wt_k)$.
\end{proof}

Tail diagnostics fitted to the observed values of $\Zh$, such as the generalised-Pareto shape estimate, cannot
detect the missing mass, because the events that carry it are absent from the sample.

\begin{corollary}[A perfect twist hides nothing \status{ours}]
If $\psi_s=h_s$, then $c_k=\wt_k\sum_jc_j$ by \eqref{eq:mismatch}, so $H_\varepsilon\le K\varepsilon$.
\end{corollary}
Hidden mass is therefore a measure of twist error. It is large when, among the particles of one run, $\psi_s$
understates $h_s$ by orders of magnitude for some of them: the frozen twist is \emph{over-confident} that the
constraint has become unreachable.

\begin{remark}[Estimating $H_\varepsilon$]
Runs with resampling switched on and off have the same distribution up to the first step at which resampling would
fire. $H_\varepsilon$ can therefore be estimated from runs without resampling: record each particle's weight at
every step, find the first step at which $\mathrm{ESS}<\tau K$, and add up the final contributions of the particles
whose normalised weight at that step was below $\varepsilon$.
\end{remark}

\begin{remark}[Why the frozen twist can be over-confident]
$\eta_t(x)$ evaluates the constraint under the denoiser's current predictions. Nothing forces those predictions to
be consistent with what the denoiser will predict later. A denoiser can assign a completion almost zero probability
at one step and a substantial probability after a few more tokens are revealed. $h_t$ averages over these later
revisions and $\eta_t$ ignores them. A synthetic denoiser whose predictions depend little on context has little
hidden mass, so tests on such models can miss the effect.
\end{remark}

\subsection{Defensive resampling}
\begin{proposition}[\status{ours}]\label{prop:defensive}
Resample from $\pi=(1-\epsilon)\wt+\epsilon/K$ and give each offspring the weight $\wt_a/\pi_a$, as in
\cref{alg:smc}. Then:
\begin{enumerate}[label=(\roman*),nosep]
  \item $\E[\Zh]=h_0$;
  \item every post-resampling weight is at most $1/(1-\epsilon)$, and every particle has at least one offspring
    with probability at least $1-(1-\epsilon/K)^K\ge1-e^{-\epsilon}$, however small its weight;
  \item for any non-negative continuation values $g_k$, the resampling step adds a conditional relative variance
    of at most $1/\epsilon$:
    \[
      \frac{\operatorname{Var}\big(\frac1K\sum_j\frac{\wt_{a_j}}{\pi_{a_j}}g_{a_j}\bigm|\mathcal F\big)}{(\sum_k\wt_kg_k)^2}
      \le\frac1\epsilon\,\frac{\sum_k(\wt_kg_k)^2}{(\sum_k\wt_kg_k)^2}\le\frac1\epsilon.
    \]
\end{enumerate}
Under multinomial resampling the same ratio is $\frac1K\big(\sum_k\wt_kg_k^2/(\sum_k\wt_kg_k)^2-1\big)$, which
has no bound when a particle with small $\wt_k$ has a large $g_k$.
\end{proposition}
\begin{proof}
(i) is \cref{thm:traj}(iii). (ii) follows from $\pi_k\ge(1-\epsilon)\wt_k$ and $\pi_k\ge\epsilon/K$. For (iii),
let $Y=(\wt_a/\pi_a)\,g_a$ with $a\sim\pi$. Then $\E Y=\sum_k\wt_kg_k$ and
$\E Y^2=\sum_k\wt_k^2g_k^2/\pi_k\le(K/\epsilon)\sum_k(\wt_kg_k)^2$. The variance of an average of $K$
independent copies is at most $\E Y^2/K$, and $\sum_k(\wt_kg_k)^2\le(\sum_k\wt_kg_k)^2$ because $g\ge0$.
\end{proof}
Taking $g_k=V_s(x^k)$ (\cref{lem:cont}) covers the situation of \cref{prop:hidden}. The construction is the
resampling counterpart of defensive importance sampling \cite{hesterberg1995}. When the twist is
accurate it costs some resampling efficiency, and in exchange its variance bound holds whatever the twist's errors.

\subsection{Tempered intermediate twists}
\begin{proposition}[\status{ours}]\label{prop:tempered}
For $\beta\in(0,1]$, let $\psi_t=\eta_t^\beta$ for $t<T$ and $\psi_T=\one_C$, with the step-exact proposal. The
potential is $G_{t+1}=\eta_t(x)^{1-\beta}\eta_{t+1}(x')^\beta/Z_d(x';x)$, and (ii) and (iii) of \cref{thm:traj}
hold. In \eqref{eq:mismatch}, $\wt_k\propto\rho^k_s\,\eta_s(x^k_s)^\beta$, so
\[
  \frac{c_k}{\wt_k}\propto h_s(x^k_s)^{1-\beta}\Big(\frac{h_s(x^k_s)}{\eta_s(x^k_s)}\Big)^\beta.
\]
\end{proposition}
\begin{proof}
$\eta_t^\beta>0$ wherever $\eta_t>0$, which includes every completable state, so \cref{prop:free} applies. The
formula for $\wt_k$ is \eqref{eq:mismatch} with $\psi_s=\eta_s^\beta$.
\end{proof}
An over-confidence factor $h_s/\eta_s=\delta^{-1}$ thus enters resampling as $\delta^{-\beta}$. The price is that
resampling uses only part of the information in the frozen twist. At $\beta\to0$, particles are resampled on the
path ratio $\rho$ alone; $\beta=1$ is TWISTER. The two remedies can be combined.

\begin{remark}[When tempering is enough]
Tempering removes hidden mass only if the tempered mismatch is large enough to be sampled. By
\cref{prop:hidden}, a particle whose normalised weight is reduced by a factor $\delta^\beta$ relative to its
share of the expectation has an offspring in some run out of $R$ only if $RK\delta^\beta\gtrsim1$, which requires
$\beta\lesssim\log(RK)/\log(1/\delta)$. When the frozen twist is over-confident by many orders of magnitude, this
forces $\beta$ close to $0$ and discards most of the look-ahead. Defensive resampling has no such condition,
because its retention probability $1-e^{-\epsilon}$ does not depend on $\delta$. It is therefore the safer default;
tempering is a complement for moderate twist errors.
\end{remark}

%======================================================================
\section{Ambiguity}\label{sec:ambiguity}
%======================================================================

By \cref{thm:traj}, an ambiguous route gives the correct target. It costs variance in two ways. The proposal
favours completions with many derivations, and the twist $\eta_d=\sum_yd(y)q(y)$ overstates how reachable $C$ is,
in proportion to the ambiguity of the remaining completions.

\begin{proposition}[Single-step cost of ambiguity \status{ours}]\label{prop:amb}
Suppose $1\le d(y)\le D$ on $C$. Importance sampling from $q_d(y)\propto d(y)q(y)\one_C(y)$ to $q(y\mid C)$ uses
weights proportional to $1/d(y)$, and its ESS fraction satisfies
\[
  \frac{(\E_{q_d}[1/d])^2}{\E_{q_d}[1/d^2]}\ge\frac{4D}{(1+D)^2}.
\]
The twist ratio satisfies $1\le\eta_d/\eta\le D$.
\end{proposition}
\begin{proof}
Let $X=1/d(y)$ with $y\sim q_d$, so $X\in[m,M]$ with $m=1/D$ and $M=1$, and write $\mu=\E X$. From
$(X-m)(M-X)\ge0$ we get $X^2\le(m+M)X-mM$, hence $\E X^2\le(m+M)\mu-mM$. The quadratic
$\frac{(m+M)^2}{4mM}\mu^2-(m+M)\mu+mM$ has discriminant zero, so it is non-negative, which gives
$\E X^2\le\frac{(m+M)^2}{4mM}\mu^2$ (Kantorovich's inequality). Substituting $m=1/D$ and $M=1$ gives the ESS
bound. The twist bound is immediate from $1\le d\le D$.
\end{proof}

The bound is informative when the ambiguity is bounded. For the RNA pairing grammar, $D$ grows exponentially with
$L$ and the bound is vacuous. Early in decoding, though, the twist only sees the ambiguity of completions consistent
with the current partial sequence, which can be much smaller than $D$. Whether a per-step measure of this kind
controls the ESS of whole trajectories is open (\cref{sec:approach}).

Exact counting that ignores derivations is out of reach for ambiguous grammars in general. For CFG constraints on
hidden Markov models, the exact constrained likelihood is NP-hard already at ambiguity 2, and an FPRAS exists when
the ambiguity is polynomially bounded \cite{marzouk2022}. This suggests approximate twists under bounded ambiguity. By \cref{prop:free}, approximate or randomised twists are allowed, provided the proposal
density is exact and the terminal twist is $\one_C$.

%======================================================================
\section{Beyond exact routes}\label{sec:beyond}
%======================================================================

\subsection{A regular superset with Boolean pruning \status{ours}}
For a CFG constraint $C$, choose a DFA for a regular superset $R\supseteq C$: the lexer or class-level DFA,
bracketing up to a depth bound with an overflow state that accepts everything, or the regular part of a schema. Use
the chain route for $R$ as the proposal, with twist
\[
  \psi_t(x)=Z_R(x;x)\cdot\mathrm{comp}_C(x),\qquad\psi_T=\one_C,
\]
where $\mathrm{comp}_C(x)\in\{0,1\}$ reports whether $x$ can be completed to an element of $C$ (an emptiness test
for the intersection of a context-free and a regular language). By \cref{prop:free} this is exact provided
\begin{enumerate}[nosep]
  \item $C\subseteq R$, so the proposal reaches every completable state, and
  \item $\mathrm{comp}_C$ never reports 0 for a completable state. Reporting 1 for a dead state only wastes a
    particle.
\end{enumerate}
Exact completability checkers satisfy condition 2. A checker that accepts a state only if one of a few sampled
fillings can be extended may reject completable states, and then biases the target. A regular \emph{subset} (an
overflow state that rejects) violates condition 1 and is biased unless the depth bound provably covers every string
of length $L$. Pruning turns one accept-or-reject decision at the end into deaths at individual steps, which
resampling can repair. $R$ should be tight enough that the per-step death rate stays well below $1-1/K$.

\paragraph{Redrawing \status{proposed}.}
Instead of letting a proposal that lands on a dead state kill its particle, redraw until $\mathrm{comp}_C$
accepts. The potential then needs the acceptance probability $a(x)$ of a single draw. Using an unbiased estimate of
$a(x)$ (the accepted fraction among $m$ extra draws) keeps the algorithm valid by the random-twist part of
\cref{prop:free}.

\subsection{Conjunctions}
Conjunctions of automaton constraints compile to product automata, whose size is the product of the factors'
sizes. A CFG intersected with a regular language is context-free, by the Bar-Hillel construction, with a grammar polynomial in both. A copy relation conjoined with local constraints stays a
factor graph as long as the combined width stays bounded. A determinant route conjoined with other constraints does
not fit any single route described here. A conjunction whose counting problem is \#P-hard is refused.

\subsection{The intractable boundary}
If a constraint family has no FPRAS, it has no polynomial-time almost-uniform sampler either \cite{jvv86}. The
trajectory layer can still be run with a superset route and a weaker twist, but its variance can grow
exponentially with $L$. Refusing such constraints is the correct behaviour for the compiler.

%======================================================================
\section{Open directions}\label{sec:approach}
%======================================================================

\begin{enumerate}[leftmargin=*]
  \item \emph{Better twists.} \Cref{prop:hidden} makes the over-confidence of the frozen twist the main practical
    limit, and \cref{prop:free} leaves the intermediate twists free. Candidates:
    \begin{enumerate}[nosep,label=(\alph*)]
      \item tempering and defensive resampling (\cref{prop:defensive,prop:tempered}), which are cheap and come with
        guarantees that do not depend on the twist;
      \item look-ahead twists, which evaluate $Z$ under the predictions made after one or more speculative
        reveals, at the cost of extra denoiser calls;
      \item mixture twists $\psi_t=(1-\lambda)\eta_t+\lambda\hat h_t$, where $\hat h_t$ is an unbiased estimate of
        $h_t$ from native rollouts and the $\eta_t$ term keeps the twist positive;
      \item learned twists $\psi_\vartheta\approx h_t$ fitted to native rollouts, as in twisted SMC for
        autoregressive models \cite{zhao2024}, with $\eta_t$ as a feature or baseline;
      \item approximate-count twists for bounded ambiguity \cite{marzouk2022}.
    \end{enumerate}
  \item \emph{Mixing the native kernel into the proposal.} $M_\alpha=(1-\alpha)M+\alpha\kappa$ has a computable
    density, and $G\le\psi_{t+1}/(\alpha\psi_t)$ for every step. It reaches regions the step-exact proposal rarely
    visits, at the price of proposals that leave $C$ and get weight zero.
  \item \emph{Trajectory ESS under ambiguity.} Bound the ESS of whole paths in terms of a per-step ambiguity
    measure, such as the spread of $d$ over completions consistent with the current partial sequence, instead of the
    global range $D$ in \cref{prop:amb}.
  \item \emph{An online guard against hidden mass.} \Cref{prop:hidden} expresses the hidden mass through
    $h_s/\psi_s$, which cannot be observed. An estimate or upper bound available during a run (for example, from a
    small population of particles that is never resampled) would turn the offline diagnostic into a check made
    while decoding.
  \item \emph{Routes for RNA and mRNA design.} A route for a fixed target structure (a tree-shaped factor graph
    over paired and unpaired positions) and charts with a bound $W$ on pair span (cost $O(LW^2|G|)$). An external
    energy model can enter through the twist, and \cref{prop:free} covers this without change.
  \item \emph{Calibrated sampling and best-of-$K$.} Exact conditional sampling should eventually beat mode-seeking
    decoding in best-of-$K$ selection as $K$ grows. Where the crossover lies is open.
\end{enumerate}

%======================================================================
\section{Summary of results}\label{sec:status}
%======================================================================
\begin{center}
\renewcommand{\arraystretch}{1.15}
\begin{tabular}{>{\raggedright}p{6.4cm}l>{\raggedright\arraybackslash}p{5cm}}
\toprule
result & status & source or note\\
\midrule
Semirings and ambiguity (\cref{prop:semiring}); MAP for any CFG & background & \cite{wang2024atlas}\\
Circuit characterisation (\cref{thm:char}) & background & \cite{broadrick2025,jvv86}\\
Spanning trees separate monotone from general circuits (\cref{rem:arb}) & background & \cite{jerrumsnir1982}\\
Step-exact proposal (\cref{lem:prop}); free twists (\cref{prop:free}) & background & \cite{twister,delmoral2004}\\
Routes with single-pass exact samplers; compiler with refusal & ours (system) & \cref{sec:routes}\\
Route-generic trajectory exactness (\cref{thm:traj}) & ours & extends \cite{twister}; compare \cite{dang2026icml}\\
Terminal-twist pitfall & ours & \\
Two-sided identity (\cref{lem:twosided}) & ours & diagnostic\\
Continuation value and hidden mass (\cref{lem:cont,prop:hidden}) & ours & \\
Defensive resampling (\cref{prop:defensive}) & ours & compare \cite{hesterberg1995}\\
Tempered intermediate twists (\cref{prop:tempered}) & ours & \\
Single-step cost of ambiguity (\cref{prop:amb}) & ours & \\
Regular superset with Boolean pruning & ours & \cref{sec:beyond}\\
Better twists; trajectory ESS; online guard & proposed & \cref{sec:approach}\\
\bottomrule
\end{tabular}
\end{center}

%======================================================================
\begin{thebibliography}{99}\small
\bibitem{broadrick2025} O.~Broadrick, S.~Agarwal, G.~Van den Broeck and M.~Bl\"aser. The limits of tractable marginalization. arXiv:2506.12020, 2025.
\bibitem{dang2026icml} M.~Dang, J.~Song, H.~Zhang, H.~Zhao, G.~Van den Broeck and S.~Ermon. Mitigating bias in locally constrained decoding via tractable proposals. ICML 2026; arXiv:2606.01926.
\bibitem{de26} M.~Dang and S.~Ermon. Constrained decoding for diffusion language models via efficient inference over finite automata. arXiv:2607.07026, 2026.
\bibitem{delmoral2004} P.~Del Moral. \emph{Feynman--Kac Formulae: Genealogical and Interacting Particle Systems with Applications}. Springer, 2004. doi:10.1007/978-1-4684-9393-1.
\bibitem{hesterberg1995} T.~Hesterberg. Weighted average importance sampling and defensive mixture distributions. \emph{Technometrics} 37(2), 1995. doi:10.1080/00401706.1995.10484303.
\bibitem{jerrumsnir1982} M.~Jerrum and M.~Snir. Some exact complexity results for straight-line computations over semirings. \emph{J.\ ACM} 29(3), 1982. doi:10.1145/322326.322341.
\bibitem{jvv86} M.~R.~Jerrum, L.~G.~Valiant and V.~V.~Vazirani. Random generation of combinatorial structures from a uniform distribution. \emph{Theor.\ Comput.\ Sci.} 43:169--188, 1986. doi:10.1016/0304-3975(86)90174-X.
\bibitem{koo2007} T.~Koo, A.~Globerson, X.~Carreras and M.~Collins. Structured prediction models via the matrix-tree theorem. In \emph{Proc.\ EMNLP-CoNLL 2007}, pp.~141--150. ACL Anthology D07-1015.
\bibitem{marzouk2022} R.~Marzouk and C.~de la Higuera. Marginal inference queries in hidden Markov models under context-free grammar constraints. arXiv:2206.12862, 2022.
\bibitem{nussinov1980} R.~Nussinov and A.~B.~Jacobson. Fast algorithm for predicting the secondary structure of single-stranded RNA. \emph{PNAS} 77(11), 1980. doi:10.1073/pnas.77.11.6309.
\bibitem{sahoo2024} S.~S.~Sahoo, M.~Arriola, Y.~Schiff, A.~Gokaslan, E.~Marroquin, J.~T.~Chiu et al. Simple and effective masked diffusion language models. arXiv:2406.07524, 2024.
\bibitem{sz26} J.~Su and W.~Zhang. Constraints are graphs, not chains: exact decoding for diffusion language models. arXiv:2609.32900, 2026.
\bibitem{twister} A.~Thimmaiah, L.~Marinov, J.~Srinivasa, H.~Vikalo, J.~J.~Li and M.~Gligoric. Twist, don't tilt: trajectory-exact constrained decoding for masked diffusion models. arXiv:2609.35609, 2026.
\bibitem{wang2024atlas} B.~Wang, D.~D.~Mau\'a, G.~Van den Broeck and Y.~Choi. A compositional atlas for algebraic circuits. arXiv:2412.05481, 2024.
\bibitem{wilson1996} D.~B.~Wilson. Generating random spanning trees more quickly than the cover time. In \emph{Proc.\ STOC 1996}. doi:10.1145/237814.237880.
\bibitem{zhao2024} S.~Zhao, R.~Brekelmans, A.~Makhzani and R.~Grosse. Probabilistic inference in language models via twisted sequential Monte Carlo. arXiv:2404.17546, 2024.
\end{thebibliography}
\end{document}
